Two EEG corpora acquired during sustained simulated-driving tasks were used. The corpora differed in their acquisition montages and data provenance, allowing the proposed method to be evaluated under different hardware setups.

The public corpus contains data from 27 participants who performed an approximately 90-minute lane-keeping task in an immersive driving simulator~\citep{cao2019multichannel}, and is referred to as Cao2018 henceforth. The Cao2018 corpus is publicly available under a Creative Commons Attribution 4.0 International licence~\citep{cao2019multichannel} and was accessed via the published figshare repository (DOI: 10.6084/m9.figshare.6427334.v5). This study used the raw recordings. All 32 recorded electrophysiological channels (30 scalp electrodes and the A1/A2 mastoid reference channels) were supplied to the pipeline. These 32 channels were Fp1, F3, Fp2, F4, FT7, FC3, C3, FC4, FT8, C4, TP7, CP3, TP8, CP4, T5, P3, O1, T6, P4, O2, F7, T3, Fz, FCz, Cz, Pz, Oz, F8, T4, CPz, A1, and A2.

The Internal corpus is a dataset collected in house, comprising data from 17 participants, recorded during a simulated lane-keeping task. The Internal corpus was collected under a protocol approved by the Human Research Ethics Committee of Universiti Pendidikan Sultan Idris (Sultan Idris Education University; reference number 2021-0290-01). All participants provided written informed consent before the recording session. EEG was acquired with a 128-channel EGI HydroCel system, with CPz as the reference. Although the original recordings used a high-density montage, the present analyses retained a predefined subset of 32 electrodes corresponding to conventional 10--20 scalp locations, following the subsampled 128-to-32-channel montage of~\citep{Hatlestad_2023_reliable}. The 32 channels are Fp1, AF3, F3, Fp2, AF4, F4, FC5, FC1, C3, FC2, FC6, C4, CP5, CP1, CP2, CP6, P8, P3, PO3, O1, PO4, P4, P7, O2, F7, T7, Fz, Cz, Pz, Oz, F8, and T8.

We divided the 32 electrodes into four anatomical regions (frontal, central, parietal, and occipital) and further divided each region into left and right hemispheres, as per~\citep{black2012}, as shown in Figure~\ref{fig:eeg_map}. In the figure, the resulting eight regions (frontal left, frontal right, central left, central right, parietal left, parietal right, occipital left, and occipital right) are abbreviated FL, FR, CL, CR, PL, PR, OL, and OR, respectively. The figure shows the scalp locations of the 32 electrodes with the blue circles represent electrodes located in both the Internal and Cao2018 montages, the yellow circles represent electrodes located only in the Internal montage, and the red circles represent electrodes located only in the Cao2018 montage.



\begin{figure}[ht]
    \centering
    \includegraphics[width=0.9\linewidth]{eeg_map}
    \caption{Scalp locations of the 32 electrodes used in this study, divided into eight anatomical regions (frontal left/right, central left/right, parietal left/right, and occipital left/right, abbreviated FL, FR, CL, CR, PL, PR, OL, and OR). Blue circles mark electrodes present in both the Internal and Cao2018 montages, yellow circles mark electrodes present only in the Internal montage, and red circles mark electrodes present only in the Cao2018 montage.}
    \label{fig:eeg_map}
\end{figure}

